\documentclass[10pt]{article}
\usepackage[T1]{fontenc}
\usepackage[utf8]{inputenc}
\usepackage[french]{babel}
\usepackage{amsmath,amssymb,mathrsfs,mathpazo}
\usepackage{geometry}
\geometry{a4paper,margin=2cm}
\usepackage{graphicx,array,tabularx,tabularray,colortbl,multirow,multicol,booktabs,enumitem,pifont,ulem,fancybox,diagbox,cancel}
\usepackage{pstricks,pstricks-add,pst-plot,pst-tree,pst-node,pst-text,pst-func,pst-math,pst-bspline,pst-3dplot,pst-eucl,pst-solides3d}
\let\pscalculate\undefined
\usepackage{pst-bezier}
\usepackage{tikz,pgfplots,tkz-tab}
\usetikzlibrary{arrows,arrows.meta,patterns,calc,positioning,trees,decorations.markings,shapes.geometric,shapes.misc,shapes.arrows,intersections,decorations.pathmorphing,decorations.pathreplacing,backgrounds,shadows,babel}
\pgfplotsset{compat=1.18}
\usepgfplotslibrary{fillbetween}
\usepackage[np]{numprint}
\usepackage{listings,makecell,pas-tableur,twoopt}
\let\barre\undefined
\usepackage[pstricks]{tablvar}
\definecolor{darkgreen}{RGB}{0,100,0}
\providecommand{\red}{\color{red}}
\providecommand{\blue}{\color{blue}}
\providecommand{\green}{\color{green}}
\providecommand{\black}{\color{black}}
\providecommand{\white}{\color{white}}
\providecommand{\gray}{\color{gray}}
\providecommand{\cyan}{\color{cyan}}
\providecommand{\magenta}{\color{magenta}}
\providecommand{\yellow}{\color{yellow}}
\providecommand{\violet}{\color{violet}}
\providecommand{\orange}{\color{orange}}
\providecommand{\euro}{\texteuro}
\providecommand{\EUR}{\texteuro}
\providecommand{\ieme}{\textsuperscript{e}}
\providecommand{\degres}{\ensuremath{^\circ}}
\providecommand{\degree}{\ensuremath{^\circ}}
\providecommand{\vv}[1]{\overrightarrow{#1}}
\providecommand{\Coord}[1]{#1}
\providecommand{\decosix}{}
\providecommand{\starredbullet}{\ensuremath{\bullet}}
\pagestyle{empty}
\setlength{\parindent}{0pt}
\usepackage{adjustbox}\DeclareUnicodeCharacter{2236}{:}\DeclareUnicodeCharacter{2190}{\ensuremath{\leftarrow}}
\providecommand{\R}{\mathbb{R}}
\providecommand{\N}{\mathbb{N}}
\providecommand{\Z}{\mathbb{Z}}
\providecommand{\Q}{\mathbb{Q}}
\providecommand{\D}{\mathbb{D}}
\providecommand{\C}{\mathbb{C}}
\providecommand{\e}{{\mathrm{e}}}
\providecommand{\vect}[1]{\overrightarrow{#1}}
\providecommand{\vectt}[1]{\overrightarrow{\mathrm{#1}}}
\providecommand{\vvt}[1]{\overrightarrow{\mathrm{#1}}}
\providecommand{\np}[2][]{{#2~#1}}
\providecommand{\vv}[1]{{\overrightarrow{#1}}}
\renewcommand{\d}{{\mathrm{d}}}
\providecommand{\degres}{{^\circ}}
\providecommand{\degre}{{^\circ}}
\providecommand{\blue}{\color{blue}}
\providecommand{\red}{\color{red}}
\providecommand{\enskip}{\quad}
\providecommand{\llbracket}{\lbrack\!\lbrack}
\providecommand{\rrbracket}{\rbrack\!\rbrack}
\providecommand{\euro}{{\text{€}}}
\providecommand{\no}{{\text{n°}}}
\providecommand{\interval}[2]{{[#1;#2]}}
\providecommand{\intervFO}[2]{{[#1;#2[}}
\providecommand{\dotfill}{\ldots}

\providecommand{\ds}{\displaystyle}
\providecommand{\textding}[1]{#1}
\providecommand{\Oij}{\(\mathrm{O};\vec{\imath},\vec{\jmath}\)}
\providecommand{\Oijk}{\(\mathrm{O};\vec{\imath},\vec{\jmath},\vec{k}\)}
\providecommand{\Ouv}{\(\mathrm{O};\vec{u},\vec{v}\)}
\providecommand{\intervFF}[2]{\ensuremath{[#1;#2]}}
\renewcommand{\euro}{\ifmmode\text{\texteuro}\else\texteuro\fi}
\DeclareUnicodeCharacter{00B0}{\ifmmode^{\circ}\else\textdegree\fi}
\DeclareUnicodeCharacter{20AC}{\ifmmode\text{\texteuro}\else\texteuro\fi}
% Énoncés : Éducation nationale. Transcription et corrigés : APMEP, Denis Vergès et contributeurs. Découpage et classement : CoopMaths, Sébastien Lozano, Éric Elter et contributeurs. Adaptation de la présentation : Maths Explorer.
% CC BY-SA : https://coopmaths.fr/www/about/
% https://coopmaths.fr/alea/static/bac/2024/tex/bac_2024_06_sujet2_asie_2.tex
\begin{document}
\noindent\textit{APMEP / CoopMaths. CC BY-SA. Présentation adaptée par Maths Explorer.}\par\medskip
\medskip

\emph{Léa passe une bonne partie de ses journées à jouer à un jeu vidéo et s'intéresse aux chances de victoire de ses prochaines parties.}

\medskip

\emph{Elle estime que si elle vient de gagner une partie, elle gagne la suivante dans 70\,\% des cas.}

\medskip

\emph{Mais si elle vient de subir une défaite, d'après elle, la probabilité qu'elle gagne la suivante est de $0,2$.}

\emph{De plus, elle pense avoir autant de chance de gagner la première partie que de la perdre.}

\emph{On s'appuiera sur les affirmations de Léa pour répondre aux questions de cet exercice.}

\medskip

Pour tout entier naturel $n$ non nul, on définit les évènements suivants: 

\begin{itemize}[label=$\bullet~$]
\item $G_n$ : \og Léa gagne la $n$-ième partie de la journée \fg{} ;
\item $D_n$ : \og Léa perd la $n$-ième partie de la journée \fg.
\end{itemize}

\smallskip

Pour tout entier naturel $n$ non nul, on note $g_n$ la probabilité de l'évènement $G_n$.

On a donc $g_1 = 0,5$.

\medskip

\begin{enumerate}
\item Quelle est la valeur de la probabilité conditionnelle $p_{G_1}\left(D_2\right)$ ?
\item Recopier et compléter l'arbre des probabilités ci-dessous qui modélise la situation pour
les deux premières parties de la journée :

\begin{center}
\pstree[treemode=R,levelsep=3.25cm]{\TR{}}
{\pstree{\TR{$G_1~$}\taput{\ldots}}
	{\TR{$G_2~$}\taput{\ldots}
	\TR{$D_2~$}\tbput{\ldots}
	}
\pstree{\TR{$D_1~$}\tbput{\ldots}}
	{\TR{$G_2~$}\taput{\ldots}
	\TR{$D_2~$}\tbput{\ldots}
	}
}
\end{center}

\item Calculer $g_2$.
\item Soit $n$ un entier naturel non nul.
	\begin{enumerate}
		\item Recopier et compléter l'arbre des probabilités ci-dessous qui modélise la situation pour les $n$-ième et $(n + 1)$-ième parties de la journée.

\begin{center}
\pstree[treemode=R,levelsep=3.25cm]{\TR{}}
{\pstree{\TR{$G_n~$}\taput{$g_n$}}
	{\TR{$G_{n+1}~$}\taput{\ldots}
	\TR{$D_{n+1}~$}\tbput{\ldots}
	}
\pstree{\TR{$D_n~$}\tbput{\ldots}}
	{\TR{$G_{n+1}~$}\taput{\ldots}
	\TR{$D_{n+1}~$}\tbput{\ldots}
	}
}
\end{center}
		\item Justifier que pour tout entier naturel $n$ non nul,
		\[g_{n+1} = 0,5g_n + 0,2.\]

	\end{enumerate}
\item Pour tout entier naturel $n$ non nul, on pose $v_n = g_n - 0,4$.
	\begin{enumerate}
		\item Montrer que la suite $\left(v_n\right)$ est géométrique.
		
On précisera son premier terme et sa raison.
		\item Montrer que, pour tout entier naturel $n$ non nul : 
		\[g_n = 0,1 \times  0,5^{n-1} + 0,4.\]
		
	\end{enumerate}
\item Étudier les variations de la suite $\left(g_n\right)$.
\item Donner, en justifiant, la limite de la suite $\left(g_n\right)$.

Interpréter le résultat dans le contexte de l'énoncé.
\item Déterminer, par le calcul, le plus petit entier $n$ tel que $g_n - 0,4 \leqslant 0,001$.
\item Recopier et compléter les lignes 4, 5 et 6 de la fonction suivante, écrite en langage Python, afin qu'elle renvoie le plus petit rang à partir duquel les termes de la suite $\left(g_n\right)$ sont tous inférieurs ou égaux à $0,4 + e$, où $e$ est un nombre réel strictement positif.













\begin{center}
\fbox{
\sffamily
\begin{tabular}{l}
1 \quad \textbf{def} seuil(e) :\\
2 \quad \qquad g = 0.5\\
3 \quad \qquad n = 1\\
4 \quad \qquad \textbf{while} \ldots  :\\
5 \quad \qquad \qquad g = 0.5 * g + 0.2 \hspace*{1cm}\\
6 \quad \qquad \qquad n = \dots\\
7 \quad \qquad \textbf{return} (n)\\
\end{tabular}
}
\end{center}
\end{enumerate}

\bigskip

\end{document}
